% Owners: Abhinav (P8.7: land classification, MCDA, learned models, sensitivity);
%         Harsh (P8.7: building similarity, suitability map and validation results)
\section{Results}\label{sec:results}

\subsection{Land classification}\label{sec:res-classification}

\paragraph{Choosing the clusters.} K-Means and Gaussian mixture models were run for $k=3$ to 10
on seven standardised inputs: the cell's shares of water, trees, crops, built and rangeland, its
mean elevation and its mean slope (Figure~\ref{fig:kscan}). K-Means separates the cells far
better at every $k$: silhouette 0.43--0.61 \citep{rousseeuw1987}, against 0.26--0.36 for the
mixture models, which also have a worse Davies--Bouldin index \citep{davies1979} at every $k$. The
silhouette of K-Means is flat from $k=5$ to $k=8$ and drops sharply at $k=9$. We use
\textbf{$\boldsymbol{k=7}$} (silhouette 0.606 on the 2023 features, 0.585 on 2018), which is the
best or within 0.003 of the best in both years. The seven clusters are the five land-cover groups
plus two sub-groups that appear in both years:
\begin{itemize}
  \item \textbf{steep natural vegetation} (mean slope \ang{13.2}, 3\,\% of the cells): the
        Bannerghatta hills, mostly ESRI rangeland;
  \item \textbf{the low western valley} (mean elevation \SI{759}{\m}, 3\,\% of the cells),
        where cropland and built land are mixed.
\end{itemize}

\begin{figure}[ht]
  \centering
  \includegraphics[width=\textwidth]{chart_k_scan}
  \caption{Choosing the number of clusters (2023 features). The dashed line marks $k=7$.}
  \label{fig:kscan}
\end{figure}

\paragraph{From clusters to classes.} Each cluster is labelled from its centroid
(Figure~\ref{fig:profiles}): water-dominated clusters are excluded; built-up where built land is
both above 30\,\% and the largest land group; forest where trees exceed 50\,\%, or where trees
and rangeland together exceed 50\,\% on slopes above \ang{8}; everything else is usable, with
cropland-dominated clusters flagged as agricultural. The second forest rule is what keeps the
Bannerghatta scrub forest out of the usable class. Finally, the exclusion mask (water, wetland,
slopes over \ang{15}, the national park) is applied cell by cell.

\begin{figure}[ht]
  \centering
  \includegraphics[width=0.95\textwidth]{chart_cluster_profiles}
  \caption{Cluster centroids (2023): mean shares, elevation and slope, coloured by z-score.
           Each row is named with the class its rule gave and its share of the cells.}
  \label{fig:profiles}
\end{figure}

\paragraph{The three-class map.} Figure~\ref{fig:class3} shows the result for 2023. Built-up land
covers \SI{348}{\km\squared}, usable land \SI{138}{\km\squared}, forest \SI{29}{\km\squared}, and
\SI{66}{\km\squared} are excluded. Between the 2018 and 2023 maps, built-up land grows by
\SI{42}{\km\squared} and usable land shrinks by \SI{44}{\km\squared}
(Figure~\ref{fig:class-areas}).

Against ESRI's own classes, collapsed to the same three, the map agrees on 96.6\,\% of the land
cells. This is expected, since the ESRI fractions are inputs, and it mainly shows that the
cluster rules don't distort them. Against the independent WorldCover map the agreement is
64.0\,\%. Almost all of the difference comes from the products themselves: WorldCover calls only
33\,\% of the area built-up, against 62\,\% in ESRI.

\begin{figure}[ht]
  \centering
  \includegraphics[width=0.7\textwidth]{map_class3_final}
  \caption{Three-class land map, 2023 (final run).}
  \label{fig:class3}
\end{figure}

\begin{figure}[ht]
  \centering
  \includegraphics[width=0.65\textwidth]{chart_class_areas}
  \caption{Class areas in the 2018 (validation) and 2023 (final) maps.}
  \label{fig:class-areas}
\end{figure}

\subsection{MCDA / AHP baseline}\label{sec:res-mcda}

The classic method uses five criteria, each rescaled to 0--1: distance to existing built-up land,
distance to a major road, slope, land cover (rangeland and bare ground easiest to develop, then
crops, then trees) and a setback from water. The AHP pairwise matrix gives built-up and road
distance 0.30 each, slope and land cover 0.16 each and water 0.09. Its consistency ratio of 0.003
is far below Saaty's 0.1 limit \citep{saaty1980}.

Tested against persistent growth from 2018 to 2022/23, the overlay reaches \textbf{AUC 0.636}
(Table~\ref{tab:auc}), below a single criterion, distance to built-up land (0.763). Scored one at
a time, only that criterion carries real signal: slope reaches 0.612, and roads, land cover and
water stay between 0.52 and 0.55, close to chance. The AHP weights give the weak criteria 70\,\%
of the weight, which dilutes the one strong signal. We deliberately did not tune the weights on
the growth label, which is the test set. A fixed, defensible expert overlay is exactly what this
baseline stands for, and it illustrates the weakness the literature describes: the weights decide
the map, and nothing checks them against outcomes \citep{malczewski2004}.

\subsection{Learned models}\label{sec:res-rf}
% Drafted by Abhinav (who built the models); Harsh to merge with the validation results.

Two random forests \citep{breiman2001, pedregosa2011} learn from 2018 surroundings only: ring
fractions at 250 and \SI{500}{\m}, distances, road density and terrain, never the cell's own land
cover. The \textbf{change RF} learns from the 355 cells that were non-built in 2018 and 2019 but
built in 2020 and 2021. The \textbf{state RF} learns built vs non-built in 2018. Hyperparameters
are tuned by spatial cross-validation with \SI{2}{\km} blocks on the training label only, and
every validation score comes from a model that never saw the cell's block.

\begin{table}[ht]
  \centering\small
  \caption{AUC on persistent growth (validation run): 12{,}777 candidate cells, 1{,}125
           of which grew by 2022/23. RF scores are out of fold.}
  \label{tab:auc}
  \begin{tabular}{@{}lc@{}}
    \toprule
    Score & AUC \\
    \midrule
    Change RF & \textbf{0.779} \\
    Distance to built-up land (baseline) & 0.763 \\
    State RF & 0.702 \\
    MCDA / AHP & 0.636 \\
    MCDA without land cover & 0.614 \\
    Random & 0.500 \\
    Building similarity, combined score & \todo{Harsh} \\
    \bottomrule
  \end{tabular}
\end{table}

The change RF is the best score so far, but only just above distance to built-up land (+0.016;
Figure~\ref{fig:roc}). Its most important input is the built-up share within \SI{500}{\m},
followed by trees and crops nearby and the distance to built-up land; terrain and water add
little (Figure~\ref{fig:importance}). Scored on its own training cells, it would claim AUC 0.997.
Most of its 2020/21 positives also grew by 2022/23, so in-sample scores would leak the answer;
out-of-fold scoring is essential. The state RF separates built from non-built cells almost
perfectly (spatial-CV AUC 0.998), but that skill does not transfer to predicting the next growth
(0.702).

\begin{figure}[ht]
  \centering
  \begin{subfigure}[t]{0.45\textwidth}
    \includegraphics[width=\textwidth]{chart_roc}
    \caption{ROC curves.}\label{fig:roc}
  \end{subfigure}\hfill
  \begin{subfigure}[t]{0.53\textwidth}
    \includegraphics[width=\textwidth]{map_score_rf_change_validation}
    \caption{Change-RF score as of 2018 (out of fold).}
  \end{subfigure}
  \caption{Predicting growth from 2018.}
\end{figure}

\begin{figure}[ht]
  \centering
  \includegraphics[width=0.9\textwidth]{chart_rf_importance}
  \caption{What the change RF relies on: permutation importance (AUC drop on held-out folds) and
           impurity importance.}
  \label{fig:importance}
\end{figure}

\paragraph{Held-out buildings.} Do surroundings identify building locations in areas a model never
saw? With reference cells (OSM buildings, plus ESRI built cells where OSM has none) held out by
\SI{2}{\km} block, a random forest ranks the held-out reference cells above the other cells with
AUC 0.984. A nearest-neighbour similarity on the same inputs reaches 0.895, and distance to
built-up alone reaches 0.887. Using both ring scales beats the \SI{500}{\m} rings alone (0.895 vs
0.871).

\subsection{Building similarity}\label{sec:res-similarity}

\todo{Harsh (P8.7): example queries A $\to$ B, aggregate similarity map and its validation}

\subsection{Suitability map and candidate sites}\label{sec:res-suitability}

\todo{Harsh (P8.7): combined score, S1--S3 classes, candidate patches, lift, TOC, LEI split, ablation}

\subsection{Sensitivity}\label{sec:res-sensitivity}

Three settings were varied (Figure~\ref{fig:sensitivity}):
\begin{itemize}
  \item \textbf{Slope threshold} (\ang{10}, \ang{15}, \ang{20}): the number of slope-excluded
        cells changes from 1{,}685 to 158, but every AUC moves by less than 0.01.
  \item \textbf{Number of clusters} ($k=5$ to 8): class shares change by about one percentage
        point and agreement with ESRI stays at 96.6--98.8\,\%. $k=5$ agrees slightly better
        with ESRI but loses the steep-vegetation cluster, so forest falls from 5.0\,\% to
        3.9\,\% of the area.
  \item \textbf{Cell size} (50, 100, \SI{200}{\m}; features and labels rebuilt each time): AUCs
        are not comparable across sizes, because the growth label itself changes with the cell,
        and at \SI{200}{\m} only 148 cells grew.
\end{itemize}
What holds in every setting is the ranking: \textbf{change RF $>$ distance to built-up $>$
MCDA}.

\begin{figure}[ht]
  \centering
  \includegraphics[width=0.95\textwidth]{chart_sensitivity}
  \caption{Temporal AUC by slope threshold and by cell size.}
  \label{fig:sensitivity}
\end{figure}
